\section{State of the art}

NYT LARGE PRINT AND 
IMX PAPER
THE FOLLOWING COMMENTED TEX IS JUST THE COPY PASTE FROM OUR PREVIOUS PAPER AT IMX

%\subsection{Digital media reading}

%\addedimx{Recent research in the interactive media experience community has begun to explore how different media-mixes, ranging from traditional articles to 360-degree video, affect user experience in immersive journalism. For instance, Bujić et al.~\cite{bujic-imx2023} compared 360-video across devices against a 'monitor-article' condition, finding that while immersive formats may increase involvement, traditional text-based articles can offer lower levels of distraction, which is critical in contemporary news consumption. This highlights the relevance of text-based media on digital platforms.}

%Within this landscape, prior research has shown that mobile devices have reshaped reading practices and have changed readers' habits~\cite{readinghabits2016, reading-periodicals2021, mobilereading2015}.
%At the same time, reading on mobile remains more demanding than on larger displays.
%Budiu~\cite{Budiu2015} identified inherent constraints—such as limited screen size, frequent interruptions, single-window views, and unstable connectivity—that require careful design considerations to support sustained reading.

%In particular, Nielsen~\cite{Nielsen2011} reported that comprehension of complex content on small displays is only about 48\% of that on desktop monitors. This gap is explained by two main factors: mobile readers see less information at once, reducing context and thus understanding, and they need to navigate through the document, which diverts attention. Moran~\cite{moran2016reading} further shows that while short, simple texts are read with similar comprehension across devices, reading speed drops significantly for longer or denser passages, highlighting the need for layout optimization and interaction strategies to improve mobile reading performance.

%These challenges become even more acute with structured, layout-driven documents. These structured documents combine dense text with visual design and spatial organization to convey meaning. \added[comment="2. Why LP?"]{Due to their distinct design elements, findings derived from linear reading cannot be generalized to layout-driven documents~\cite{ozretic_dosen_key_2018,hollander_e-reader_2011}.} They pose distinctive challenges due to their complex structures, often organized around visual entry points~\cite{holsanova_entry_2006}. 

%added{It is worth noting that, in theory, web-responsive adaptation principles could address the general challenge of reading complex documents on small screens. However, preserving spatial semantics is often at odds with standard web-adaptation paradigms. Techniques such as responsive design handle small displays primarily through linear re-flow mechanisms~\cite{reflow, marcotte_responsive_2010}, collapsing multi-column grids into a single vertical stack defined by the DOM order. While this method preserves sequential reading order and eliminates bi-directional scrolling, it sacrifices the 2D spatial paradigm that defines newspapers and other layout-driven documents~\cite{chesham_master_2003, chiou_2024}.}

%\added{As a result, for documents whose navigation, visual hierarchy, and interpretation rely on spatial organization~\cite{holmqvist_role_2005, ohara_comparison_1997}, standard responsive design—and similarly summary-based or linearized representations—cannot serve as viable alternatives: they remove visual entry points, suppress layout cues, and fundamentally transform the nature of the reading task.
%}

%This limitation motivates evaluating large-print layouts as an alternative that preserves spatial organization while restoring entry-point legibility, allowing us to examine \replacedimx{the behavioral and performance cost of manipulable navigation (\condonenot) compared to direct structural access (\condtwonot).}{the trade-off between manipulable control (\condonenot) and global awareness (\condtwonot)}.


%\subsection{Situational visual impairments (SVIs) in mobile reading}

%Situational visual impairments~\cite{svi-challenges} (SVIs) arise when reading is disrupted by device limitations, such as small font or display size and low resolution, or by environmental factors, including luminosity and ambient noise. Tigwell et al.~\cite{tigwell-2018} investigated how SVIs like glare or movement affect mobile content usability, highlighting that designers often lack specialized tools and guidelines to address these challenges. They proposed preliminary recommendations and emphasized the need for improved design support through better resources and educational frameworks. Building on this, another study~\cite{svi-bright-env} specifically examined the impact of bright outdoor lighting on screen readability. Their results indicate that, while ambient brightness reduces readability, the intrinsic brightness and contrast of content play an even larger role, suggesting that design interventions should prioritize optimizing content contrast over solely compensating for glare.

%In parallel, technical solutions have been proposed to mitigate SVIs. \emph{SituFont}~\cite{yue2024situfont} is a dynamic font adaptation system that adjusts attributes such as size, weight, and spacing based on real-time sensor inputs and human-in-the-loop feedback. 
%The authors validated their approach through comparative evaluations across eight simulated SVI scenarios, considering factors such as environment (indoor vs. outdoor), user mobility, and luminosity. Their findings show that SituFont significantly improves reading efficiency and reduces cognitive and physical workload compared to manual adjustments, highlighting the potential of adaptive font systems for supporting mobile reading under situational constraints. 

%While these studies provide valuable insights and technical solutions for reading under SVIs, they primarily focus on text consumption. Our work extends this line of research by investigating structured digital layout-based documents \replacedimx{under constrained visual access scenarios (\CVAS), examining the specific behavioral and performance cost of magnification strategies.}{, where reading is affected under constrained visual access scenarios (\CVAS).} In particular, we examine SVIs related to small headline font sizes and the need for magnification.

%\subsection{Reading with low vision}

%\addedimx{Despite the importance of inclusive design, users with disabilities have historically been underrepresented within the IMX community. A survey of 17 years of research~\cite{vatavu-imx2021} revealed that only 4.23\% of papers addressed accessibility, with a specific call for more empirical studies involving people with disabilities. This gap motivates our work, which evaluates reading behavior across two distinct scenarios of visual limitation. }

%In this context, digital devices provide new opportunities for readers \LV to access information through adaptable text formats and presentation options~\cite{reading-digital-legge}. A growing body of research has examined how these readers interact with digital reading environments, including studies that simulate \CVAS to better understand performance under magnification and limited visual access. In this topic, a key contribution by Atilgan et al.~\cite{atilgan2020} provides a unified framework to evaluate how print size and display size impact reading speed for both small-display readers and individuals with low vision using magnified text. Their results show that limitations in font and display size can prevent some readers from maximizing reading performance. 

%Several studies have examined display configurations and strategies for magnified reading. Xiong et al.~\cite{xiong_digital_2022} proposed guidelines emphasizing the importance of exceeding the critical print size (CPS) and maintaining at least 13 characters per line to support fluent reading. Granquist et al.~\cite{granquist_how_2018} asked participants with low vision to adjust text for comfortable reading, finding that most relied on enlarging text, and to a lesser extent on reducing viewing distance. Beckmann et al.~\cite{beckmann_psychophysics_1996} studied manual navigation with CCTVs (stand-mounted video magnifiers) and reported that reading speed decreases unless the viewing window is sufficiently wide (>10 characters), consistent with the threshold identified by Xiong et al.~\cite{xiong_digital_2022}. Magnification and contrast enhancement were also shown to improve reading speed in simulated low-vision environments, with magnification generally having a more substantial effect than contrast~\cite{magnification-low-vision}.

%Other work has examined navigation and magnification strategies in more detail. Bowers et al.\cite{bowers-reading} analyzed reading with optical magnifiers, highlighting challenges such as line retrace that need to be addressed in future research. Tang et al.~\cite{tang_screen_2023} compared two types of screen magnification on modern devices: full-screen magnification (analogous to pan-and-zoom) and lens magnification (allowing users to see a larger portion of the un-magnified screen for spatial reference). They identified trade-offs in performance and usability, showing that lens mode led to more consistent and uniform mouse movements, whereas full mode caused longer and more frequent pauses. Similarly, Aguilar and Castet~\cite{aguilar-castet2017} developed a gaze-controlled system that magnifies a portion of text while maintaining a global view, which participants found more comfortable than traditional CCTVs.

%Together, these studies highlight the impact of magnification on low-vision reading. However, most prior work focuses on plain, unstructured text, analyzing reading speed, errors, and navigation patterns within a fixed set of lines. In contrast, our study investigates structured layout-based documents, such as newspapers, where reading involves navigating not only text but also the spatial layout. We examine how visual access limitations affect document reading for both individuals with \NV and those \LV. In particular, we examine the behavioral and performance cost of gesture-based magnification (\condonenot) compared to the direct structural access of large-print layouts (\condtwonot), quantifying how each strategy impacts the reading process. 